
\chapter{Bias sui Bias}

\begin{quotation}
``Il bias più pericoloso è quello di credere di non averne''
\end{quotation}

Esiste una dimensione meta-cognitiva dei bias che è particolarmente insidiosa: i bias sui bias stessi. Non solo sistematicamente sbagliamo nei nostri giudizi, ma sbagliamo anche nel valutare quando e quanto stiamo sbagliando.

\section{Il Punto Cieco dei Bias}

La maggior parte delle persone riesce facilmente a riconoscere i bias cognitivi negli altri, ma ha una ``blind spot'' quando si tratta dei propri. È come avere un punto cieco nella vista: sappiamo che esiste, ma non riusciamo a vederlo direttamente.

Questo ``bias del punto cieco'' è particolarmente forte nelle persone più istruite, che tendono a pensare: ``Io ho studiato psicologia, so come funzionano questi meccanismi, quindi sono immune''. È l'equivalente cognitivo di pensare che conoscere la teoria della guida ti renda automaticamente immune agli incidenti stradali.

Emily Pronin di Princeton lo ha dimostrato in modo brillante\footnote{Pronin, E., Lin, D. Y., \& Ross, L. (2002). ``The bias blind spot: Perceptions of bias in self versus others''. \textit{Personality and Social Psychology Bulletin}, 28(3), 369-381.}. Ha chiesto a studenti di valutare quanto loro e i loro pari fossero influenzati da vari bias. Risultato prevedibile: tutti pensavano di essere meno biased degli altri. Ma ecco il colpo di genio: anche dopo aver spiegato loro il bias del punto cieco, continuavano a pensare di esserne meno affetti degli altri. È un bias che sopravvive alla propria diagnosi.

\section{L'Effetto Dunning-Kruger del Dunning-Kruger}

Conoscete l'effetto Dunning-Kruger, vero? Quella curva dove i meno competenti sopravvalutano le proprie capacità mentre gli esperti le sottovalutano? Bene, c'è un meta-livello ancora più perverso: le persone che hanno appena imparato cos'è l'effetto Dunning-Kruger pensano immediatamente di esserne immuni.

È quello che chiamo ``l'effetto Dunning-Kruger del Dunning-Kruger'': la conoscenza superficiale di un bias ci rende paradossalmente più vulnerabili ad esso, perché ci dà una falsa sensazione di immunità. È come prendere una dose omeopatica di vaccino e pensare di essere invincibili.

Ho visto questo fenomeno ripetersi infinite volte nelle conferenze di psicologia. Qualcuno presenta l'effetto Dunning-Kruger, il pubblico ride di quelli ``stupidi che non sanno di essere stupidi'', e nessuno ,NESSUNO ,si ferma a pensare: ``Aspetta, e se fossi io quello sulla parte sinistra della curva?''

È il narcisismo intellettuale nella sua forma più pura: usiamo la conoscenza dei bias come prova che non ne siamo affetti. È come usare il fatto di possedere uno specchio come prova di non avere niente tra i denti.

\section{Il Bias di Superiorità del Bias}

C'è un fenomeno ancora più sottile che ho osservato tra chi studia scienze cognitive: quello che chiamo il ``bias di superiorità del bias''. Funziona così: impari a riconoscere i bias cognitivi, inizi a vederli ovunque negli altri, e sviluppi un senso di superiorità intellettuale che è esso stesso un bias.

``Guarda quei poveri idioti che credono alle fake news per il bias di conferma.''
``Ovvio che ha investito in crypto, classica overconfidence.''
``La sua relazione è fallita per il bias del costo affondato.''

Diventiamo esperti nel diagnosticare i bias altrui mentre siamo completamente ciechi ai nostri. È come un oculista miope che prescrive occhiali a tutti tranne che a se stesso.

Il problema è che questa ``competenza diagnostica'' ci dà l'illusione di essere osservatori neutrali della follia umana, quando in realtà siamo partecipanti attivi. È la versione cognitiva del guardarsi un documentario sulla natura pensando di essere David Attenborough invece che uno degli animali ripresi.

\section{La Trappola della Metacognizione}

Più diventiamo consapevoli dei nostri processi mentali, più rischiamo di cadere in quello che chiamo ``la trappola della metacognizione''. Pensiamo che osservare i nostri pensieri ci renda immuni dai loro difetti. È come pensare che guardare il proprio stomaco digerire ci renda immuni dall'indigestione.

La metacognizione ,pensare sui propri pensieri ,è certamente utile. Ma porta con sé il suo set di bias. Il ``bias dell'introspezione'', per esempio: la convinzione che abbiamo accesso privilegiato e accurato ai nostri processi mentali. Spoiler: non ce l'abbiamo.

Nisbett e Wilson lo dimostrarono già nel 1977\footnote{Nisbett, R. E., \& Wilson, T. D. (1977). ``Telling more than we can know: Verbal reports on mental processes''. \textit{Psychological Review}, 84(3), 231-259.}: quando chiedi alle persone perché hanno preso certa decisione, inventano spiegazioni plausibili che hanno poco a che fare con i veri meccanismi causali. Ma sono convintissime che quelle spiegazioni siano accurate.

È come chiedere al vostro fegato perché ha processato l'alcol in un certo modo. Non lo sa. E voi non sapete perché avete scelto quella marca di cereali al supermercato, anche se potete inventare una storia convincente al riguardo.

\section{Il Paradosso dell'Educazione ai Bias}

Ecco una verità scomoda: più educhiamo le persone sui bias cognitivi, più creiamo nuove forme di bias. È quello che succede in ogni corso di pensiero critico, workshop sulla razionalità, o libro divulgativo sui bias (sì, incluso questo).

Le persone imparano a riconoscere una manciata di bias famosi ,conferma, ancoraggio, disponibilità ,e pensano di aver scaricato l'antivirus mentale definitivo. In realtà hanno solo imparato a riconoscere alcuni virus mentre ne contraggono di nuovi.

È particolarmente evidente su Twitter e Reddit, dove orde di neo-esperti di bias cognitivi si accusano a vicenda di fallacia logiche con la stessa sophistication di bambini che giocano a ``chi ce l'ha più lungo'' con i nomi dei dinosauri. ``Il tuo argomento è una fallacia del piano inclinato!'' ``No, il TUO è un ad hominem!'' ``Quello che hai appena detto è un esempio perfetto di bias di conferma!''

Usano la conoscenza dei bias come arma retorica invece che come strumento di autoanalisi. È come usare il manuale di psichiatria per diagnosticare malattie mentali a tutti tranne che a se stessi.

\section{L'Illusione dell'Obiettività}

Forse il meta-bias più pericoloso è l'illusione dell'obiettività: più conosciamo i bias, più ci convinciamo di poter raggiungere una prospettiva neutrale e oggettiva. È il sogno impossibile dell'osservatore imparziale.

I giornalisti lo chiamano ``view from nowhere'' ,la vista da nessun luogo. I filosofi lo chiamano ``the God's eye view''. Gli psicologi lo chiamano ``delusione''.

Non esiste una prospettiva libera da bias. Anche il tentativo di essere obiettivi è filtrato attraverso i nostri bias su cosa significhi essere obiettivi. È bias all the way down, come tartarughe cognitive impilate all'infinito.

Il fisico Niels Bohr una volta disse: ``Il contrario di una verità banale è un errore. Il contrario di una verità profonda può essere un'altra verità profonda.'' Lo stesso vale per i bias: il contrario di un bias non è l'obiettività, è spesso solo un bias diverso.

\section{Il Bias del Cinismo}

C'è una tentazione, arrivati a questo punto, di cadere in quello che chiamo ``il bias del cinismo'': siccome tutti hanno bias e nessuno può essere veramente obiettivo, allora tutto è ugualmente sbagliato, tutte le opinioni sono ugualmente invalide, e possiamo anche rinunciare a cercare di pensare meglio.

È il nichilismo epistemologico, la resa cognitiva totale. ``Siamo tutti biased, quindi chi se ne frega.'' È attraente nella sua semplicità, consolante nella sua universalità. Se tutti sono stupidi, nessuno lo è davvero.

Ma questo è, ovviamente, un altro bias. Il bias di credere che siccome la perfezione è impossibile, il miglioramento è inutile. È come dire che siccome non possiamo eliminare completamente le malattie, tanto vale abolire la medicina.

\section{La Recursione Infinita}

A questo punto potreste pensare: ``Ok, quindi ho bias sui miei bias, e probabilmente ho bias sui miei bias sui miei bias, e...'' Sì, esatto. È recursione infinita, specchi che riflettono specchi, un loop cognitivo senza fine.

Ogni livello di consapevolezza porta con sé i propri punti ciechi. Ogni tentativo di correggere un bias ne introduce di nuovi. È come il principio di indeterminazione di Heisenberg applicato alla mente: non puoi osservare i tuoi processi mentali senza modificarli.

Ma ecco il paradosso finale: questa consapevolezza della recursione infinita è essa stessa soggetta a bias. Potreste sopravvalutare quanto è profonda (``Wow, sono così meta!'') o sottovalutare quanto influenza il vostro pensiero (``Ok, lo so, quindi sono a posto'').

\section{L'Umiltà Epistemica (Che È Anch'essa un Bias)}

Allora, cosa ci resta? Forse quello che i filosofi chiamano ``umiltà epistemica'': riconoscere i limiti della propria conoscenza. Ma attenzione: anche l'umiltà epistemica può diventare un bias, una posa intellettuale, un modo per sentirsi superiori a chi è ``arrogantemente sicuro''.

``Io so di non sapere'' diceva Socrate. Ma anche questa è una forma di sapere, e porta con sé la sua forma di arroganza. È l'arroganza dell'umiltà, il bias di non avere bias, l'orgoglio di riconoscere i propri limiti.

È un gioco che non si può vincere. Ogni mossa per uscire dal labirinto dei bias è essa stessa una mossa dentro il labirinto. È come cercare di saltare sopra la propria ombra: più in alto salti, più lunga diventa.

\section{Il Ponte verso il Pianeta che Bolle}

E questo ci porta perfettamente al prossimo capitolo, dove vedremo come tutti questi meta-bias si manifestano nel dibattito sul cambiamento climatico. Perché se c'è un argomento dove i bias sui bias raggiungono livelli da capogiro, è proprio quello.

Da una parte, chi nega il cambiamento climatico accusa gli altri di bias di conferma, di seguire il gregge scientifico, di catastrofismo motivato politicamente. Dall'altra, chi accetta il consenso scientifico accusa i negazionisti di bias di conferma, di seguire il gregge politico, di ottimismo motivato economicamente.

Entrambi i gruppi sono convinti di essere quelli razionali, quelli che vedono attraverso i bias degli altri. Entrambi usano la conoscenza dei bias come arma per delegittimare l'opposizione. Entrambi sono certi di essere immuni dai bias che diagnosticano negli altri.

È il teatro dell'assurdo dei meta-bias: tutti esperti nel riconoscere i bias altrui, tutti ciechi ai propri, tutti convinti di essere l'eccezione. E intanto il pianeta si scalda, indifferente alle nostre sofisticate analisi metacognitive.

Ma di questo parleremo nel prossimo capitolo. Per ora, vi lascio con questa chicca di saggezza paradossale: il primo passo per ridurre i propri bias è accettare che non potrete mai eliminarli completamente. Il che, ovviamente, potrebbe essere solo il mio bias che parla.

Il loop continua. I bias si moltiplicano. E noi, cavernicoli con gli smartphone che studiano se stessi, continuiamo a inciampare sui nostri piedi mentre analizziamo l'andatura.

Benvenuti nella sala degli specchi cognitiva. L'uscita è un'illusione. Ma almeno il viaggio è interessante.
